Therefore, from the previous observations, it has been decided to design another type of SAT solver, named \textbf{Conflict-Driven Clause Learning}. Whenever 
a contradiction happens, it is set a process composed by the following operations:
\begin{enumerate}
    \item \textbf{Clause learning}.
    \begin{itemize}
        \renewcommand{\labelitemi}{-}
        \item Augments the original CNF $X$ propositional formula with an additional \textbf{conflict clause} that summarizes the root cause of the contradiction.
        \item From the point above, we will learn from the mistakes done and ignore huge sections of the search space that would be useless to explore.
    \end{itemize}
    \item \textbf{Backjumping}.
    \begin{itemize}
        \renewcommand{\labelitemi}{-}
        \item Based on the cause of the contradiction, skips the decisions and derivations that did not lead to conflict, going higher in the search tree.
    \end{itemize}
\end{enumerate}
\begin{example}
    e.g. Dumb example. \vspace{7pt}

    Starting from the following list of clauses:
    \begin{enumerate}
        \item $\neg p \vee q$
        \item $\neg r \vee s$
        \item $\neg t \vee \neg u$
        \item $\neg q \vee \neg t \vee u$
    \end{enumerate} \vspace{3.5pt}
    Define if the whole CNF propositional formula is satisfiable or unsatisfiable. \vspace{7pt}

    \textbf{Resolution}. \vspace{7pt}
    \begin{center}
        \begin{tabular}{|p{2.75cm}|p{5cm}|}          
            \hline
                \textbf{Rule} & $M$\\ 
            \hline
                Decide & $p^d$ \\
            \hline
                UnitPropagate & $p^d \ q$ \\
            \hline
                Decide & $p^d \ q \ r^d$ \\
            \hline
                UnitPropagate & $p^d \ q \ r^d \ s$ \\
            \hline
                Decide & $p^d \ q \ r^d \ s \ t^d$  \\
            \hline
                UnitPropagate & $p^d \ q \ r^d \ s \ t^d \ u$  \\
            \hline
                & \\
            \hline
        \end{tabular}
    \end{center} \vspace{7pt}
    Observing what we did so far, we have reached a contradiction. In particular, the literals $t$ and $u$ are set to true but, however, the third clause
    ($\neg t \vee \neg u$) requires at least one of them assigned to false. We backtrack to the last decision done $t^d$ and setting the literal $t$ to 
    false ($\neg t$) we state the satisfiability of the problem. \vspace{7pt}

    Anyway, supposing we have a bigger CNF formula, composed by more than four clauses. Even though we define the literal $t$ equal to false ($\neg t$), 
    we cannot achieve the satisfiability of the problem: we backtrack to the last decision before $t^d$, namely $r^d$, setting its value to another 
    truth value. Later, we will try again the decision $t^d$ failing in the same conflicts as before. \vspace{7pt}
    
    Briefly, even if we backtrack changing truth values, we end always in the same contradictions, failing exactly in the same way. \vspace{7pt}

    It would be nice to understand something more by a contradiction analysis.
\end{example}
For the contradiction between the literals $t$, $u$ and the third clause ($\neg t \vee \neg u$), the decision $r^d$ has nothing to do with them. What we 
really want to understand are the clauses that contribute to the contradiction described in the example: literal $u$ was derived using the fourth clause 
($\neg q \vee \neg t \vee u$) and the contradiction was found using the third clause. \vspace{7pt}

For sure we fail for the third clause ($\neg t \vee \neg u$), but also we reach a contradiction combining it with the fourth one (\textit{it sets the boolean 
variable $u$ to true. Therefore, there is a kind of relationship between the third and the fourth clauses}). The main idea is to apply directly the resolution 
process to these clauses getting a new \textbf{conflict clause}, obtained as follows:
$$ 5. \ \neg q \vee \neg t $$
The clause $5.$ is something we learnt after the failure, taking this new information we can do unit propagation in a different way. This newly conflict clause 
gives us a clue on where to \textbf{backjump}.
\begin{example}
    e.g. Extending the same previous example. \vspace{7pt}

    \textbf{Resolution}. \vspace{7pt}
    \begin{center}
        \begin{tabular}{|p{2.75cm}|p{5cm}|}          
            \hline
                \textbf{Rule} & $M$\\ 
            \hline
                Decide & $p^d$ \\
            \hline
                UnitPropagate & $p^d \ q$ \\
            \hline
                Decide & $p^d \ q \ r^d$ \\
            \hline
                UnitPropagate & $p^d \ q \ r^d \ s$ \\
            \hline
                Decide & $p^d \ q \ r^d \ s \ t^d$  \\
            \hline
                UnitPropagate & $p^d \ q \ r^d \ s \ t^d \ u$  \\
            \hline
                Backjumping & $p^d \ q \ \neg t$ \\
            \hline
        \end{tabular}
    \end{center} \vspace{7pt}
\end{example}